\chapter{Offener Messdatenvertrag und Suchbefunde}
\index{Messdaten}\index{Kalibrierung}
\section{Stand der Fallannahme}
Die maßgeblichen Quellen sind \path{Doc/65_VALIDATION_DATA_CONTRACT.md} und
\path{Doc/66_M9_DATA_SEARCH_AND_RELEASE_OPTIONS.md}. Weder die ursprüngliche
noch die erneute öffentliche Suche hat einen vollständigen unabhängigen
quasistatischen Piezofall zur Annahme geführt. Issues 63 und 65 bleiben offen
außerhalb des neu geplanten M9. Ihre ursprünglichen Anforderungen wurden
nicht durch den synthetischen Unsicherheitsfall erfüllt.

Die erneute Sichtung vom 10. Oktober 2026 prüfte Repository-Metadaten,
zwei kleine CSVs, eine TXT-Datei und die Blatt-/Textstruktur eines XLSX-Archivs.
Sie führte keine fremden Auswertungsskripte aus, digitalisierte keine Kurven,
fittete keine Materialwerte und berechnete keinen Messvergleich. Eine
zugängliche Datei mit Datenlizenz ist nur ein Teil des Annahmevertrags.

\section{Kandidaten und konkrete Ablehnungsgrenzen}
\subsection*{Green Piezo, Zenodo 7920056 v1.}
CC BY 4.0 und Design und Charakterisierungsdaten sind nachgewiesen.
Das 776-MB-Archiv wurde nicht vollständig auditiert; unabhängige vollständige
Tensoren, Unsicherheiten und quasistatische Messschaltung bleiben unqualifiziert.
\subsection*{Light-Annealed Films, Zenodo 15526822.}
Die gesichtete Feld-/Verschiebungsdatei ist unter CC BY 4.0 und nennt 100 Hz.
Das belegt allein keinen quasistatischen Zustand. Der aus derselben
Verschiebung abgeleitete effektive Koeffizient wäre keine unabhängige Eingabe
für die Prüfung genau dieser Antwort.
\subsection*{Polaron-defect complexes, Zenodo 21454715 v2.}
CC BY 4.0; reale räumliche Verschiebungs- sowie Dateien für Bias und Frequenz.
Die beschriebene DC-aktivierte elektrostriktive Kilohertz-Antwort besitzt
keinen angenommenen linearen quasistatischen Probenvertrag.
\subsection*{Relaxor-Kristalle, Figshare 19346039 v1.}
CC BY 4.0 und 14 Abbildungsblätter mit Materialcharakterisierung.
Die gesichtete Struktur belegt keinen vollständigen unabhängigen statischen
Gerätevergleich mit allen erforderlichen Metadaten.
\subsection*{Quarz-Stimmgabel, arXiv:1203.5767v1.}
Ein wissenschaftlich interessanter statischer Kalibrierfall mit Unsicherheit.
Ein separat lizenziertes Rohdatenarchiv und der vollständige unabhängige
Geometrie, Elektroden, Kristallachsen und Tensoren wurden nicht angenommen.
\subsection*{Duval et al., Southampton 263266.}
Messungen an PZT-Mikrocantilevern sind publiziert. Der Repository-Volltext ist
auf Mitarbeiter beschränkt; zugängliche Rohdateien und Rechte fehlen im Nachweis.
\subsection*{Gesing, Figshare 5766708.}
Die Repositorybeschreibung nennt CC BY 4.0 und einen frequenzabhängigen
Beschleunigungssensor. Das ist kein statischer d33-Verschiebungsfall;
Rohdateizugriff wurde in dieser Recherche nicht erfolgreich geprüft.
Die sechs früheren Kandidaten bleiben mit ihren damaligen Prüfgrenzen im
Messdatenvertrag erhalten. Numerische und DFT-Daten werden trotz offenen
Downloads nicht als experimentelle Messungen bezeichnet.

\section{Herkunft und Nutzungsrechte}
Die Evidenzliste \path{Doc/evidence/issue63/research-2026-10-10.json} enthält
URLs, Versionen, Rechte und SHA-256 tatsächlich gesichteter Dateien.
Die kleinen unveränderten Textdateien im Repository tragen Attribution;
das XLSX wird referenziert. Ausschnitt-Snapshots der APIs besitzen eigene
Hashes, getrennt vom Hash der vollständigen abgerufenen Antwort.
Keiner dieser Bestände ist ein eingefrorener angenommener Validierungsfall.

Dieses Handbuch enthält eigene Abbildungen aus eigenen synthetischen
Auswertungen und eigene Benchmarkrohwerte unter der Projektlizenz MIT.
Fremde Messdateien werden im Handbuch nicht erneut abgebildet oder als
eigene Werte verarbeitet. Die Primärlinks und Dateiattribution stehen im
versionierten Suchbericht. Verfügbarkeit einer Publikation ersetzt keine
gesonderte Prüfung der Datenrechte.

\section{Anforderungen an einen späteren Vergleich}
Vor einer Rechnung sind Geometrie und tolerierte Maße, Materialtensoren und
deren unabhängige Herkunft, SI-Einheiten, Voigt-/Achsenkonventionen,
Einspannung, Elektroden, tatsächlich anliegende Spannung und Messschaltung
festzulegen. Sensorposition und Messoperator müssen zur FEM-Zielgröße passen:
ein flächengemittelter Weg ist kein Knotenmaximum.

Kalibrier- und Validierungsläufe sind getrennt einzufrieren. Materialparameter
aus der zu prüfenden Zielantwort nachzufitten und dieselbe Antwort anschließend
als Validierung zu verwenden, verletzt die Unabhängigkeit. Ein späterer
Holdout muss vor der Anpassung gewählt und von ihr unberührt bleiben.

Für eine Differenz $\Delta=F_{\mathrm{FEM}}-F_{\mathrm{Messung}}$ sind
Numerik und Messunsicherheit, Eingangskovarianzen und mögliche gemeinsame
Fehlerquellen zu dokumentieren. Nur unter begründeten Unabhängigkeitsannahmen
dürfen passende Varianzen summiert werden. Eine zulässige Abweichung braucht
einen vorab definierten Verwendungszweck und Messvertrag; das Handbuch erfindet
keine universelle Akzeptanzzahl. Ein vollständiger zukünftiger Nachweis benötigt
ein positives Datengate und die separate experimentelle Studie.
